\section{Preliminaries}\label{sec:prelim}
\subsection{Sequential Design and Information Loss}
\label{sec:background_pipeline}
Accelerator design commonly proceeds from a resource budget to a hardware
template, then to operation mappings and a memory schedule. Engineers iterate
these choices using analytical models, simulation, and compiler feedback.
Automated flows replace individual choices with search tools, but an early
commitment still limits what downstream stages can recover. A mapper that
exports only its fastest isolated tile discards alternatives that might reduce
workspace, enable fusion, or retain data for a later consumer. Joint design
keeps those alternatives available until their complete resource costs are known.

\subsection{Linear Relaxation and Recovery}
\label{sec:lps}
An LP minimizes a linear objective subject to linear rows and coordinate
bounds. Discrete implementation selections are relaxed to $[0,1]$; bounded
products are represented by convex envelopes. The resulting optimum is a
lower bound only for the domain actually covered by the formulation and its
certificate. Fractional selections need not describe a chip. Recovery chooses
integral hardware and implementations, then checks dependencies, storage,
transfers, and predicted latency. A feasible recovered design is useful even
when the remaining optimality interval is wide.

\subsection{Capability Comparison}
\label{sec:capability-matrix}
Table~\ref{tab:capability-matrix} compares decision spaces, rather than claiming
that every listed capability has a measured gain. CHEETAH's row describes the
master formulation; Section~\ref{sec:decode_decomp} identifies the subset exercised
by the current full-inference execution. Related Work explains restricted cells.
